\documentclass[10pt,a4paper]{amsart}
\usepackage[margin=19mm]{geometry}
\usepackage{amsmath,amssymb,mathtools}
\usepackage[scr=rsfs,cal=euler]{mathalfa}
\usepackage{xfrac}
\usepackage[unicode,hidelinks,bookmarks=false]{hyperref}
\newcommand{\ie}{i.e.\ }
\newcommand{\cf}{cf.\ }
\newcommand{\resp}{respectively\ }
\newcommand{\Subsection}{\subsection}
\providecommand{\nobreakdash}{\nobreak\hbox{-}}
\setlength{\emergencystretch}{2em}
\multlinegap0pt
\usepackage{bm}
\usepackage{bbm}
\usepackage{dsfont}
\usepackage{xspace}
\usepackage{cancel}
\usepackage{mathtools}
\usepackage{amsthm}
\makeatletter
\@ifundefined{theorem}{\newtheorem{theorem}{Theorem}}{}
\@ifundefined{proposition}{\newtheorem{proposition}[theorem]{Proposition}}{}
\makeatother
\NewDocumentCommand{\conn}{o}{\nabla\IfValueT{#1}{^{#1}}}
\newcommand{\Sc}{\mathrm{Sc}}
\newcommand{\bnd}[1]{\partial #1}
\newcommand{\cbar}[1]{\dom{c}\left(#1\right)}
\newcommand{\dom}[1]{\hat{#1}}
\newcommand{\inner}[1]{\langle #1 \rangle}
\newcommand{\onorm}[1]{\Vert #1 \Vert}
\newcommand{\pnorm}[1]{| #1 |}
\title[A scalar-mean curvature comparison theorem for manifolds wit]{A scalar-mean curvature comparison theorem for manifolds with iterated conical singularities}
\author{Milan Jovanovic, Jinmin Wang}
\date{Source: arXiv:2506.24059v1 (CC BY 4.0)}
\begin{document}
\maketitle

\noindent\emph{Source and licence.} A scalar-mean curvature comparison theorem for manifolds with iterated conical singularities. Version arXiv:2506.24059v1 (CC BY 4.0), distributed under the Creative Commons Attribution 4.0 International licence, as recorded in the arXiv licence metadata for this exact version and shown on the abstract page https://arxiv.org/abs/2506.24059v1. Full rights notice: licenses/arxiv-2506.24059-CC-BY-4.0.txt.\par\smallskip

\noindent\textit{Editorial scope.} Counted unit: the Proposition whose source label is \texttt{prop:scalar-mean curvature estimate} in the article (source0.tex lines 756--770, source label \texttt{prop:scalar-mean curvature estimate}), delivered together with the article's own complete original proof of it (source0.tex lines 771--856, one \texttt{proof} environment of 86 lines). No mathematics is written, repaired or re-derived: statement and proof are the article's own text line for line. Required context retained uncounted: thm:scalar curvature estimate (lines 644--651). Every definition, display and label of those lines is retained unchanged. Omitted from this package and not redistributed: 1--643, 652--755, 857--1986. Nothing inside a retained excerpt is omitted or altered: every retained body is delivered exactly as the article's lines are written, so no omission receipt of any kind accompanies this package. Citations: the retained excerpts carry no citation command at all, so no bibliography entry of the article is transcribed and no reference is redistributed. Reference adaptation: none was needed and none was made; every reference inside the delivered unit resolves inside it, so no unresolved reference or citation is produced. Printed slips kept literal: no printed word, symbol, hypothesis or number of the counted statement or of its proof was altered. Layout and class: the article's own class is \texttt{amsart} with packages including amssymb, mathrsfs, latexsym, amsmath, amsthm, amsfonts, dsfont, upgreek, mathtools, epsfig, cite, amscd, graphicx, tikz-cd; the wrapper is the lane's pinned \texttt{amsart} editorial wrapper at 10pt on A4, which restates the theorem-like environments the retained text needs and adds the article's own macro definitions that the retained text needs (\texttt{\textbackslash Sc}, \texttt{\textbackslash bnd}, \texttt{\textbackslash cbar}, \texttt{\textbackslash conn}, \texttt{\textbackslash dom}, \texttt{\textbackslash inner}, \texttt{\textbackslash onorm}, \texttt{\textbackslash pnorm}), with extra wrapper packages (bm, bbm, dsfont, xspace, cancel, mathtools). The delivered numbering is the unit's own. Assets: no retained span contains a figure environment, a graphics-inclusion command, a PDF-page inclusion, a table or a TikZ picture; no image file is copied into this package, the package declares no asset dependency and no image file is redistributed with it. Licence: version arXiv:2506.24059v1 of this article is distributed under the Creative Commons Attribution 4.0 International licence, as recorded in the arXiv licence metadata for this exact version and shown on the abstract page https://arxiv.org/abs/2506.24059v1. A full rights notice is delivered at licenses/arxiv-2506.24059-CC-BY-4.0.txt. Characters: every retained excerpt is pure ASCII, so the delivered engine, XeLaTeX with its default Latin Modern text font, reports no missing character.
\begin{theorem}\label{thm:scalar curvature estimate}
    Let \((\dom{M}, \bnd{\dom{M}}, \dom{g})\) and \((M, \bnd{M}, g)\) be manifolds with FACS and \(F\) an asymptotically conical map between them.
    If the curvature operator \(\mathcal{R}\) on M is nonnegative and if \(\sigma\) is a smooth section of \(S(T\dom{M} \oplus F^*TM)\) vanishing near the boundary and near the singular stratum of \(\dom{M}\), then the following inequality holds:
    \begin{equation}\label{eqn:scalar curvature estimate}
        \int_{\dom{M}}\pnorm{D\sigma}^2\geq \int_{\dom{M}}\pnorm{\nabla \sigma}^2+
        \int_{\dom{M}} \frac{1}{4}(\dom{\Sc}-\onorm{\wedge^2dF}F^*\Sc )\pnorm{\sigma}^2.
    \end{equation}
\end{theorem}

\begin{proposition}\label{prop:scalar-mean curvature estimate}
    Let \((\dom{M}, \bnd{\dom{M}}, \dom{g})\) and \((M, \bnd{M}, g)\) be manifolds with FACS and smooth boundaries \(\bnd{\dom{M}}\), \(\bnd{M}\).
    Let \(F\colon \dom{M} \to M\) be an asymptotically conical map and let \(F_\partial\) denote its restriction to the boundary.
    Suppose the curvature operator \(\mathcal{R}\) on M is nonnegative and suppose the second fundamental form \(\mathcal{A}\) along the boundary of \(\bnd{M}\) is also nonnegative.
    If \(\sigma\) is a smooth section of \(S(T\dom{M} \oplus F^*TM)\) vanishing near the singular stratum of \(\dom{M}\), then the following inequality holds:
    \begin{multline}
        \int_{\dom{M}}\pnorm{D\sigma}^2\geq \int_{\dom{M}}\pnorm{\nabla \sigma}^2+\int_{\dom{M}} \frac{1}{4}(\dom{\Sc}-\onorm{\wedge^2dF}F^*\Sc )\pnorm{\sigma}^2\\
        + \int_{\bnd{\dom{M}}} \inner{D^{\partial, \partial} \sigma, \sigma} + \int_{\bnd{\dom{M}}} \frac{1}{2} (\dom{H} - \onorm{d(F_\partial)} (F_\partial)^*H )\pnorm{\sigma}^2.
    \end{multline}
    Here, \(D\) is the twisted Dirac operator on \(S(T\dom{M} \oplus F^*TM)\)  and \(D^{\partial, \partial}\) is the boundary Dirac operator defined as
    \begin{equation*}
        D^{\partial, \partial} \sigma = \sum_i \cbar{\dom{\nu}}\cbar{\dom{e}_i} \conn[\partial, \partial] \sigma.
    \end{equation*}
    In the above, \(\dom{\nu}\) is the inner unit normal and the vectors \(\{\dom{e}_i\}\) form a local orthonormal frame of the boundary of \(\dom{M}\).
\end{proposition}

\begin{proof}
    Observe that
    \begin{equation*}
        \int_{\dom{M}} \pnorm{D \sigma}^2 = \int_{\dom{M}} \inner{D^2 \sigma, \sigma} + \int_{\bnd{\dom{M}}} \inner{\cbar{\dom{\nu}} D \sigma, \sigma}.
    \end{equation*}

    Rewriting the last term in the above
    \begin{align*}
        \int_{\bnd{\dom{M}}} \inner{\cbar{\dom{\nu}} D \sigma, \sigma} &= -\int_{\bnd{\dom{M}}} \inner{\conn[]_{\dom{\nu}} \sigma, \sigma} + \int_{\bnd{\dom{M}}} \inner{\sum_i \cbar{\dom{\nu}} \cbar{\dom{e}_i} \conn[]_{\dom{e}_i} \sigma, \sigma}\\
        &= -\int_{\bnd{\dom{M}}} \inner{\conn[]_{\dom{\nu}} \sigma, \sigma} + \int_{\bnd{\dom{M}}} \inner{\sum_i \cbar{\dom{\nu}} \cbar{\dom{e}_i} \conn[\partial, \partial] \sigma, \sigma}\\
        &\qquad\qquad + \int_{\bnd{\dom{M}}} \inner{\frac{1}{2} \sum_{i, j} \inner{\conn[]_{\dom{e}_i} \dom{\nu}, \dom{e}_j} \cbar{\dom{\nu}} \cbar{\dom{e}_i} \cbar{\dom{\nu}} \cbar{\dom{e}_j} \sigma, \sigma}\\
        &\qquad\qquad + \int_{\bnd{\dom{M}}} \inner{\frac{1}{2} \sum_{i, j} \inner{\conn[]_{F_*\dom{e}_i} \nu, e_j} \cbar{\dom{\nu}} \cbar{\dom{e}_i} c(\nu) c(e_j) \sigma, \sigma}.
    \end{align*}
    To simplify the notation, define \(\dom{c}_\partial(\dom{e}_i) \coloneqq \dom{c}(\dom{\nu}) \dom{c}(\dom{e}_i)\) and \(c_\partial(e_i) \coloneqq c(\nu) c(e_i)\).
    Rewriting the above,
    \begin{equation}\label{eqn:boundary dirac operator term}
        \begin{split}
            \int_{\bnd{\dom{M}}} \inner{\cbar{\dom{\nu}} D \sigma, \sigma} &= -\int_{\bnd{\dom{M}}} \inner{\conn[]_{\dom{\nu}} \sigma, \sigma} + \sum_i \int_{\bnd{\dom{M}}} \inner{D^{\partial, \partial} \sigma, \sigma}\\
            &\qquad\qquad + \int_{\bnd{\dom{M}}} \inner{\underbrace{\frac{1}{2} \sum_{i, j} \inner{\conn[]_{\dom{e}_i} \dom{\nu}, \dom{e}_j} \dom{c}_\partial(\dom{e}_i) \dom{c}_\partial(\dom{e}_j)}_{\text{Term 1}}\sigma , \sigma}\\
            &\qquad\qquad + \int_{\bnd{\dom{M}}} \inner{\underbrace{\frac{1}{2} \sum_{i, j} \inner{\conn[]_{F_*\dom{e}_i} \nu, e_j} \dom{c}_\partial(\dom{e}_i) c_\partial(e_j)}_{\text{Term 2}}\sigma , \sigma}.
        \end{split}
    \end{equation}

    Isolating Term 1, we compute
    \begin{equation*}
        \frac{1}{2} \sum_{i, j} \inner{\conn[]_{\dom{e}_i} \dom{\nu}, \dom{e}_j} \dom{c}_\partial(\dom{e}_i) \dom{c}_\partial(\dom{e}_j) = \frac{1}{2} \sum_i \inner{\dom{\nu}, \conn[]_{\dom{e}_i} \dom{e}_i} - \underbrace{\frac{1}{2} \sum_{i \neq j} \inner{\dom{\nu}, \conn[]_{\dom{e}_i} \dom{e}_j} \cbar{\dom{e}_i} \cbar{\dom{e}_j}}_{=\, 0}
    \end{equation*}
    where the last term is 0 since \(\inner{\dom{\nu}, \conn[]_{\dom{e}_i} \dom{e}_j}\) is symmetric with respect to \(\dom{e}_i\) and \(\dom{e}_j\) while \(\cbar{\dom{e}_i} \cbar{\dom{e}_j}\) is skew-symmetric with respect to \(\dom{e}_i\) and \(\dom{e}_j\) when \(i \neq j\).
    Hence, Term 1 is just the mean curvature of \(\bnd{\dom{M}}\).

    Next, we analyze Term 2.
    Start by taking frames \(\{\dom{e}_i\}\) and \(\{e_i\}\) for which \(F_*\dom{e}_i = \mu_i e_i\).
    Observe that \(\inner{\conn[]_{F_*\dom{e}_i} \nu, e_j} = -\inner{\mathcal{A} F_*\dom{e}_i, e_j}\).
    If \(B\) denotes the square root of \(\mathcal{A}\), then
    \begin{align*}
        -\inner{\mathcal{A} F_*\dom{e}_i, e_j} &= -\inner{B F_*\dom{e}_i, B e_j} = -\sum_k \inner{BF_*\dom{e}_i, e_k} \inner{Be_j, e_k}
    \end{align*}

    Hence, Term 2 can be rewritten as
    \begin{align*}
        \frac{1}{2} \sum_{i, j} \inner{\conn[]_{F_*\dom{e}_i} \nu, e_j} \dom{c}_\partial(\dom{e}_i) c_\partial(e_j) &= -\frac{1}{2}\sum_{i,j,k} \inner{BF_*\dom{e}_i, e_k} \inner{Be_j,e_k} \dom{c}_\partial(\dom{e}_i) c_\partial(e_j) \\
        &= -\frac{1}{2} \sum_k \dom{c}_\partial(\dom{B} e_k) c_\partial(B e_k) \\
        &= \frac{1}{4} \sum_k \Big(\alpha^{-2} \dom{c}_\partial(\dom{B}e_k)^2 + \alpha^2 c_\partial(Be_k)^2 \\
        &\qquad\qquad\qquad -\big(\alpha^{-1} \dom{c}_\partial(\dom{B}e_k) + \alpha c_\partial(Be_k)\big)^2\Big)
    \end{align*}
    where \(\dom{B}e_k \coloneqq \sum_i \left< F_*\dom{e}_i, Be_k \right> \dom{e}_i\) and \(\alpha\) is some scalar function to be specified.
    Notice that \(\alpha^{-1} \dom{c}_\partial(\dom{B} e_k) + \alpha c_\partial(B e_k)\) is skew-symmetric and hence the negative of its square is nonnegative.
    At the same time,
    \begin{align*}
        \frac{1}{4} \alpha^2 \sum_k c_\partial(Be_k)^2 = -\frac{1}{4} \alpha^2 \sum_k \pnorm{B e_k}^2 &= -\frac{1}{4} \alpha^2 \sum_k \left<\mathcal{A}e_k, e_k\right> \\
        &= -\frac{1}{4} \alpha^2 F^*H.
    \end{align*}
    and also
    \begin{align*}
        \frac{1}{4} \alpha^{-2} \sum_k \dom{c}_\partial(\dom{B} e_k)^2 &=  \frac{1}{4} \alpha^{-2} \sum_{k,i,j} \inner{F_*\dom{e}_i, Be_k} \inner{F_*\dom{e}_j, Be_k} \dom{c}_\partial(\dom{e}_i) \dom{c}_\partial(\dom{e}_j)\\
        &= -\frac{1}{4} \alpha^{-2} \sum_{k,i} \left<F_*\dom{e}_i, Be_k\right>^2 \\
        &= -\frac{1}{4} \alpha^{-2} \sum_{k,i} \left<\mu_i Be_i, e_k\right>^2 \\
        &\geq -\frac{1}{4} \alpha^{-2} \sum_i \onorm{dF}^2 \pnorm{B e_i}^2 \\
        &= -\frac{1}{4} \alpha^{-2} \onorm{dF}^2 F^*H.
    \end{align*}

    We have established the following estimate for Term 2:
    \begin{align*}
        \frac{1}{2} \sum_{i, j} \inner{\conn[]_{F_*\dom{e}_i} \nu, e_j} \dom{c}_\partial(\dom{e}_i) c_\partial(e_j) &\geq -\frac{1}{4} \alpha^{-2} \onorm{dF}^2 F^*H - \frac{1}{4} \alpha^2 F^*H \\
        &= -\frac{\onorm{dF}}{2} F^*H
    \end{align*}
    by taking \(\alpha = \sqrt{\onorm{dF}}\).

    To summarize,
    \begin{multline}\label{eqn:final estimate of dirac boundary term}
        \int_{\bnd{\dom{M}}} \inner{\cbar{\dom{\nu}} D \sigma, \sigma} \geq -\int_{\bnd{\dom{M}}} \inner{\conn[]_{\dom{\nu}} \sigma, \sigma} + \int_{\bnd{\dom{M}}} \inner{D^{\partial, \partial}\sigma, \sigma}\\
        + \int_{\bnd{\dom{M}}} \frac{\dom{H}}{2} \pnorm{\sigma}^2 - \int_{\bnd{\dom{M}}} \onorm{dF} \frac{F^*H}{2} \pnorm{\sigma}^2.
    \end{multline}

    At the same time, by the proof of Theorem \ref{thm:scalar curvature estimate},
    \begin{equation}\label{eqn:final estimate of square of dirac operator term}
        \int_{\dom{M}} \inner{D^2 \sigma, \sigma} = \int_{\dom{M}} \inner{\nabla^*\nabla \sigma, \sigma} + \int_{\dom{M}} \frac{\dom{\Sc}}{4} \pnorm{\sigma}^2 - \int_{\dom{M}} \frac{F^*\Sc}{4} \pnorm{\sigma}^2.
    \end{equation}.

    Integrating by parts,
    \begin{equation}\label{eqn:integration by parts of laplacian}
        \int_{\dom{M}} \inner{\nabla^*\nabla \sigma, \sigma} = \int_{\dom{M}} \pnorm{\conn[] \sigma}^2 + \int_{\bnd{\dom{M}}} \inner{\conn[]_{\dom{\nu}} \sigma, \sigma}.
    \end{equation}

    Combining Equations \eqref{eqn:final estimate of dirac boundary term}, \eqref{eqn:final estimate of square of dirac operator term}, and \eqref{eqn:integration by parts of laplacian} yields the estimate.
\end{proof}

\end{document}
